\section{Controlled User Study}
\label{sec:user-study}

We conducted a controlled study to examine whether behavioral patch review
helps workflow owners understand and govern a skill repair relative to a
capability-matched conversational interface. The comparison concerns how
repair work is represented and reviewed, rather than the presence of any
single interface control. We evaluated the complete interaction model because
its contribution lies in connecting intent, artifact revision, and enacted
behavior; the study does not estimate the isolated effect of each workspace
component.

\subsection{Research Questions and Conditions}

The study addressed three research questions:

\begin{description}[leftmargin=0pt,labelindent=0pt,style=nextline]
    \item[\textbf{RQ1: Behavioral understanding.}] Does the structured
    workspace help owners understand the enacted behavioral change produced by
    a candidate revision?
    \item[\textbf{RQ2: Repair scope.}] Does it help owners construct and
    maintain a repair scope across intended changes, valued behavior, and
    unresolved boundaries?
    \item[\textbf{RQ3: Release judgment.}] Does it support decisions that are
    better calibrated to the available evidence while preserving the owner's
    authority over release?
\end{description}

In the \emph{workspace} condition, behavioral commitments, the skill
revision, and matched execution evidence remain persistently visible and
cross-linked. In the \emph{chat} condition, participants can access the same
skill, cases, evidence-generating operations, and revision capabilities
through conversation, but these objects are presented sequentially rather
than integrated into a persistent review representation. Both conditions ask
participants to make the same substantive judgments about the motivating
incident, neighboring cases, candidate revision, and final disposition.

The two conditions use the same underlying model, candidate-generation
process, typed tool runtime, scenario worlds, case exposure, outcome criteria,
and execution budgets. Interface organization is therefore the primary
manipulation: the workspace externalizes relationships among intent, artifact,
and enacted behavior, whereas the chat condition requires participants to
maintain and reconcile those relationships through the conversation. Exact
configurations and execution protocols are reported in
Appendix~\ref{app:system-details}.

\subsection{Design, Tasks, and Assignment}

We used a counterbalanced within-participant design with two domains:
travel-disruption rebooking and expense-receipt review. Each participant
completed one task in each condition and encountered each domain only once,
avoiding direct repetition of the same repair problem. Four assignment
sequences balanced condition order, domain order, and condition--domain
mapping. We determined the final sample size through
\DataNeeded{DESCRIBE THE A PRIORI POWER ANALYSIS, TARGET EFFECT, POWER, AND
ALLOWANCE FOR EXCLUSIONS} after a usability pilot whose participants did not
enter confirmatory analyses.

Each task placed the participant in the role of a workflow owner responsible
for reviewing a problematic execution and deciding whether a revised skill is
ready for reuse. The scenario includes the motivating incident, visible
contrastive cases that probe desired changes and preservation constraints, a
policy boundary, and cases reserved for post-decision assessment. Selected
cases were withheld from candidate generation to distinguish direct repair of
visible examples from transfer to behavior that did not guide the revision.
Both conditions disclosed that all actions occurred in an isolated
environment, but neither revealed the injected fault, a reference repair, or
the hidden assessment criteria.

\subsection{Participants}

We recruited \StudyN{} adults who
\DataNeeded{REPORT ELIGIBILITY CRITERIA, INCLUDING EXPERIENCE WITH LLM TOOLS,
RECURRING WORKFLOWS, OR REVIEW OF AUTOMATED OUTPUTS}. We did not require
participants to program or directly author an agent skill because the target
role was the person responsible for acceptable behavior and release. The
sample included \DataNeeded{REPORT AGE SUMMARY, GENDER ONLY AS COLLECTED,
EDUCATION OR OCCUPATIONAL BACKGROUND, LLM EXPERIENCE, AND RELEVANT DOMAIN
EXPERIENCE}. \DataNeeded{REPORT WHETHER AND HOW MANY PARTICIPANTS HAD PRIOR
EXPERIENCE EDITING PROMPTS, SKILLS, OR AUTOMATIONS}.

Participants received \StudyCompensation{} for a session lasting
\DataNeeded{PLANNED AND OBSERVED SESSION DURATION}. The study was reviewed
under \StudyEthics{}. We excluded \DataNeeded{NUMBER} participants according
to the preregistered criteria because \DataNeeded{REPORT EXCLUSION REASONS};
\DataNeeded{NUMBER} completed participant records entered the confirmatory
analysis. A separate \DataNeeded{PILOT N} participants completed the usability
pilot and were not included in the reported results.

\subsection{Procedure}

Each task followed four phases shared across conditions. First, participants
inspected the skill and a recurring execution failure, then described what
should have happened. Second, they examined contrastive cases and recorded
whether the revision should change the observed behavior, preserve it, or
leave the policy boundary unresolved. These judgments established a
prospective repair scope before candidate outcomes were available.

Third, participants reviewed bounded evidence connecting the current behavior to
the source skill and inspected a candidate revision produced from the same
information in both conditions. They then compared the complete original and
revised skills across the reviewed cases. Participants could inspect
underlying executions, edit or regenerate the candidate, request further
evidence, revise their scope, defer the repair, or release the candidate. The
interface recorded scope changes made after outcomes became visible so that
they could be distinguished from pre-reveal commitments.

Finally, before a previously unseen case was executed, participants predicted
observable aspects of the candidate's behavior and reported their confidence.
They then completed experience and workload measures. The final artifact was
evaluated on the prediction case and on a research holdout that was unavailable
during the task. These post-decision cases assessed behavioral understanding
and transfer without changing the participant's completed release decision.

\subsection{Measures}

We preregistered three co-primary outcomes aligned with the research questions
\DataNeeded{PROVIDE PREREGISTRATION REFERENCE AND STATE ANY DEVIATIONS}.
\emph{Held-out behavioral understanding} combines accuracy on observable
predictions about a previously unseen case with calibration derived from
participants' confidence. \emph{Behavioral-scope adherence} measures the
proportion of pre-reveal Change and Preserve commitments satisfied by the final
candidate; cases intentionally left unresolved do not enter this denominator.
\emph{Evidence-calibrated decision quality} assesses whether the terminal
action and its rationale are supported by the evidence available at the time.
A release is unsupported when an unmet change target, preservation regression,
execution failure, or material evidence gap remains; a decision not to release
is evidence-grounded when it identifies a remaining mismatch, unresolved
consequence, or evidentiary limitation. A change within an unresolved boundary
requires acknowledgement but is not treated as an error by default.

Secondary outcomes characterized both repair quality and interaction process.
They included incident resolution, preservation of visible and generator-
withheld behavior, performance on the blind holdout, identification of case
roles, regression detection, task time, evidence requests, candidate and scope
revisions, release path, and consequential failures. Post-task items measure
participants' understanding of the original and revised skills, ability to
anticipate behavior, awareness of repair boundaries, perceived sufficiency of
evidence, explainability of the final decision, and perceived control. We also
collected the six Raw TLX dimensions and report them separately.

\subsection{Analysis and Data Integrity}

Confirmatory models included condition, domain, period, and order as fixed
effects and participant as a random intercept. We used mixed-effects models
appropriate to binary, continuous, and count outcomes, applied the
preregistered multiplicity procedure to the three co-primary outcomes, and we
report effect sizes with confidence intervals. The two domains were treated as
fixed task contexts, not as a random sample of all possible agent workflows.

The system recorded the skill and scope versions, timing of outcome reveals,
candidate provenance, execution traces, evidence requests, and terminal
decision for each task. This record supports reconstruction of the evidence
available to a participant without exposing reference repairs or hidden
holdouts during the interaction. Questionnaire responses were joined to this
record only after the task was complete.

Before confirmatory data collection, we completed technical acceptance,
calibrated the faulty and reference behaviors across all study cases, tested
the end-to-end candidate path, audited both conditions for information
leakage, and conducted a separate usability pilot. We then froze the study
build; \DataNeeded{REPORT ANY POST-FREEZE DEVIATIONS OR STATE THAT NONE
OCCURRED}. Appendix~\ref{app:system-details} documents the exact configuration
and reconstruction protocol.

\subsection{Results}

\begin{table*}[t]
    \caption{Co-primary outcomes by condition. Values and model estimates are
    placeholders pending the confirmatory analysis. Higher values indicate
    better performance for prediction accuracy and scope adherence; lower
    Brier scores indicate better calibrated predictions.}
    \label{tab:user-study-primary}
    \small
    \begin{tabularx}{\textwidth}{@{}p{0.25\textwidth}Y Y Y@{}}
        \toprule
        \textbf{Outcome} & \textbf{Workspace} & \textbf{Chat} &
        \textbf{Condition effect} \\
        \midrule
        Held-out prediction accuracy &
        \DataNeeded{ESTIMATE AND CI} &
        \DataNeeded{ESTIMATE AND CI} &
        \DataNeeded{MODEL EFFECT, CI, ADJUSTED P} \\
        Prediction calibration (Brier) &
        \DataNeeded{ESTIMATE AND CI} &
        \DataNeeded{ESTIMATE AND CI} &
        \DataNeeded{MODEL EFFECT, CI, ADJUSTED P} \\
        Behavioral-scope adherence &
        \DataNeeded{ESTIMATE AND CI} &
        \DataNeeded{ESTIMATE AND CI} &
        \DataNeeded{MODEL EFFECT, CI, ADJUSTED P} \\
        Evidence-calibrated decision quality &
        \DataNeeded{ESTIMATE AND CI} &
        \DataNeeded{ESTIMATE AND CI} &
        \DataNeeded{MODEL EFFECT, CI, ADJUSTED P} \\
        \bottomrule
    \end{tabularx}
\end{table*}

\paragraph{Task completion and condition integrity.}
Participants completed \DataNeeded{NUMBER OF ANALYZED TASKS} analyzed tasks.
\DataNeeded{REPORT INTERRUPTIONS, RESUMED TASKS, EXECUTION FAILURES, AND ANY
MISSING QUESTIONNAIRES}. Audit records confirmed that the two conditions used
the same scenario artifacts, candidate-generation inputs, execution evidence,
and outcome criteria in \DataNeeded{PROPORTION OF TASKS; EXPLAIN ANY
DEVIATIONS}. The manipulation check showed that participants experienced the
workspace as a more persistent and connected representation of the repair
\DataNeeded{REPORT ITEM OR CHECK, EFFECT, CI, AND P}, while access to the skill,
cases, and repair operations did not differ substantively
\DataNeeded{REPORT CAPABILITY-PARITY CHECK}. We found
\DataNeeded{REPORT PERIOD, ORDER, AND DOMAIN EFFECTS, INCLUDING ANY
CONDITION--DOMAIN INTERACTION}.

\paragraph{RQ1: The workspace improved understanding of enacted behavioral
change.}
Participants predicted behavior on the previously unseen case more accurately
in the workspace condition than in chat (Table~\ref{tab:user-study-primary}).
The estimated difference was \DataNeeded{PREDICTION-ACCURACY EFFECT, CI,
ADJUSTED P, AND EFFECT SIZE}. Their confidence was also better calibrated:
selected-answer Brier scores were lower by \DataNeeded{BRIER EFFECT, CI,
ADJUSTED P}. \DataNeeded{REPORT WHETHER THE BENEFIT CAME FROM CORRECTLY
ANTICIPATING DECISIONS, TOOL ACTIONS, CONFIRMATION ORDER, OUTCOMES, OR A
COMBINATION}.

The behavioral prediction result was consistent with participants' experience
ratings. They reported greater understanding of the candidate
\DataNeeded{ITEM ESTIMATE, EFFECT, CI, P} and greater ability to anticipate its
behavior \DataNeeded{ITEM ESTIMATE, EFFECT, CI, P}. Interaction records suggest
that workspace participants used the cross-links among cases, commitments, and
the skill diff to reconstruct why behavior changed
\DataNeeded{REPORT RELEVANT NAVIGATION OR EVIDENCE-VIEWING PATTERN}. In chat,
participants more often requested or produced repeated summaries of earlier
cases \DataNeeded{COUNT OR PROPORTION AND COMPARISON}, indicating that the same
evidence was harder to maintain as a single review object. A participant
described the contrast as \emph{``\DataNeeded{VERIFIED POST-TASK QUOTE ABOUT
CONNECTING THE REVISION TO BEHAVIOR ACROSS CASES}.''}

\paragraph{RQ2: The workspace produced repairs that better respected their
prospective scope.}
Final candidates in the workspace condition satisfied a greater proportion of
the participants' pre-reveal Change and Preserve commitments
\DataNeeded{SCOPE-ADHERENCE EFFECT, CI, ADJUSTED P, AND EFFECT SIZE}. Incident
success was \DataNeeded{WORKSPACE AND CHAT RATES}, while preservation of
visible behavior was \DataNeeded{WORKSPACE AND CHAT RATES}. The difference was
especially apparent on the generator-withheld Preserve case
\DataNeeded{TRANSFER EFFECT, CI, P}, which did not directly inform candidate
generation. Blind-holdout performance showed \DataNeeded{REPORT DIRECTION,
MAGNITUDE, CI, AND WHETHER THE RESULT WAS CONFIRMATORY OR EXPLORATORY}.

Participants using the workspace identified more preservation regressions
before their terminal decision \DataNeeded{EFFECT, CI, P} and were more likely
to acknowledge a material change in an Unresolved case
\DataNeeded{EFFECT, CI, P}. They made \DataNeeded{WORKSPACE VERSUS CHAT SCOPE
REVISIONS} scope revisions after seeing candidate behavior. Inspection of
these revisions showed that \DataNeeded{REPORT HOW OFTEN THEY CLARIFIED A
GENUINE POLICY BOUNDARY VERSUS RETROFITTED INTENT TO THE CANDIDATE}. This
distinction is important: behavioral patch review supported revising an
underspecified scope, but did not treat post-reveal agreement as evidence that
the original repair had been well scoped.

\paragraph{RQ3: Release decisions were better calibrated to available
evidence.}
Workspace participants reached an evidence-supported terminal decision in
\DataNeeded{WORKSPACE RATE} of tasks, compared with \DataNeeded{CHAT RATE} in
chat \DataNeeded{MODEL EFFECT, CI, ADJUSTED P}. They released a candidate while
a recorded Change or Preserve mismatch remained in \DataNeeded{WORKSPACE RATE}
versus \DataNeeded{CHAT RATE} of eligible tasks, and they explicitly cited an
unresolved consequence or evidence limitation when declining release in
\DataNeeded{WORKSPACE RATE} versus \DataNeeded{CHAT RATE}. The result was not
simply greater reluctance to act: when the candidate satisfied the visible
scope and evidence was stable, workspace participants released it in
\DataNeeded{RATE}, compared with \DataNeeded{RATE} in chat.

Decision rationales in the workspace condition more often referred to both a
desired change and a preservation or boundary consideration
\DataNeeded{CODING PROCEDURE, AGREEMENT, RATES, EFFECT, AND P}. Participants
also reported greater evidence sufficiency \DataNeeded{EFFECT, CI, P},
decision explainability \DataNeeded{EFFECT, CI, P}, and perceived control
\DataNeeded{EFFECT, CI, P}. One participant summarized the division of labor:
\emph{``\DataNeeded{VERIFIED QUOTE ABOUT THE SYSTEM DOING EVIDENCE OR
IMPLEMENTATION WORK WHILE THE PARTICIPANT RETAINED THE RELEASE DECISION}.''}

\paragraph{Workload and interaction cost.}
The improvement in the co-primary outcomes did not require reliably longer
tasks: completion time differed by \DataNeeded{TIME EFFECT, CI, P}, with median
times of \DataNeeded{WORKSPACE MEDIAN} and \DataNeeded{CHAT MEDIAN}. Workspace
participants inspected \DataNeeded{MORE, FEWER, OR SIMILAR} distinct evidence
items and requested \DataNeeded{MODEL OR EXECUTION CALL COMPARISON}. Raw TLX
mental demand was \DataNeeded{DIRECTION, EFFECT, CI, P}; the remaining Raw TLX
dimensions showed \DataNeeded{REPORT PATTERN WITHOUT COLLAPSING DIMENSIONS INTO
AN UNJUSTIFIED COMPOSITE}. We observed \DataNeeded{NUMBER AND DESCRIPTION} tasks
in which the workspace's additional structure slowed an otherwise simple
repair, providing a boundary on when the interaction is beneficial.

\subsection{Study Summary}

Across two task domains, the results were consistent across the three
levels of behavioral patch review. The structured workspace improved
participants' ability to anticipate the enacted delta, helped their final
revision adhere to a prospectively recorded behavioral scope, and supported
terminal decisions that better reflected the available evidence. Process and
experience measures suggest that the benefit came from keeping cases,
commitments, skill changes, and executions available as one reviewable object,
rather than from providing additional agent capabilities. The study evaluates
the integrated interaction model and does not identify which individual visual
representation is necessary or sufficient for the observed effects.
